\documentclass{article}
\usepackage{mathtools} 
\usepackage{fontspec}
\usepackage[UTF8]{ctex}
\usepackage{amsthm}
\usepackage{mdframed}
\usepackage{xcolor}
\usepackage{amssymb}
\usepackage{amsmath}


% 定义新的带灰色背景的说明环境 zremark
\newmdtheoremenv[
  backgroundcolor=gray!10,
  % 边框与背景一致，边框线会消失
  linecolor=gray!10
]{zremark}{说明}

% 通用矩阵命令: \flexmatrix{矩阵名}{元素符号}{行数}{列数}
\newcommand{\flexmatrix}[4]{
  \[
  #1 = \begin{pmatrix}
    #2_{11}     & #2_{12}     & \cdots & #2_{1#4}   \\
    #2_{21}     & #2_{22}     & \cdots & #2_{2#4}   \\
    \vdots      & \vdots      & \ddots & \vdots     \\
    #2_{#31}    & #2_{#32}    & \cdots & #2_{#3#4}
  \end{pmatrix}
  \]
}

% 简化版命令（默认矩阵名为A，元素符号为a）: \quickmatrix{行数}{列数}
\newcommand{\quickmatrix}[2]{\flexmatrix{A}{a}{#1}{#2}}

\begin{document}
\title{注释 6.4}
\author{张志聪}
\maketitle

\begin{zremark}
  定义 6.13 的证明
\end{zremark}

\textbf{证明：}

由条件，可得$f$在$x_0$处的n阶泰勒公式：
\begin{align*}
  \lim\limits_{x \to x_0} \frac{f(x) - [f(x_0) + f'(x_0)(x - x_0) + \frac{f''(x_0)}{2!}(x - x_0)^2 + \cdots + \frac{f^{(n)}(x_0)}{n!}(x - x_0)^n]}{(x - x_0)^n} = 0
\end{align*}
因为$f^{(k)}(x_0) = 0 (k = 1, 2, \cdots, n-1)$，所以上式等于：
\begin{align*}
  \lim\limits_{x \to x_0} \frac{f(x) - [f(x_0) + \frac{f^{(n)}(x_0)}{n!}(x - x_0)^n]}{(x - x_0)^n} = 0 \\
  \lim\limits_{x \to x_0} \frac{f(x) - f(x_0)}{(x - x_0)^n} - \frac{f^{(n)}(x_0)}{n!} = 0 \\
\end{align*}
所以 
\begin{align*}
  \lim\limits_{x \to x_0} \frac{f(x) - f(x_0)}{(x - x_0)^n} = \frac{f^{(n)}(x_0)}{n!}
\end{align*}

(i) $n$是偶数时，$(x - x_0)^n > 0$，当$f^{(n)}(x_0) < 0$，由极限的保号性可知，
存在$U(x_0)$使得
\begin{align*}
  f(x) - f(x_0) < 0
\end{align*}
所以，$f$在$x_0$处取极大值；

类似地，当$f^{(n)}(x_0) > 0$，$f$在$x_0$处取极小值。

(ii) $n$为奇数时，$x > x_0$时，$(x - x_0)^n > 0$，当$f^{(n)}(x_0) < 0$，由极限的保号性可知，
存在$U_+(x_0)$使得
\begin{align*}
  f(x) - f(x_0) < 0 \\
  f(x) < f(x_0)
\end{align*}

当$x < x_0$时，$(x - x_0)^n < 0$，当$f^{(n)}(x_0) < 0$，由极限的保号性可知，
存在$U_+(x_0)$使得
\begin{align*}
  f(x) - f(x_0) > 0 \\
  f(x) > f(x_0)
\end{align*}

所以，$f$在$x_0$处不取极值。

同理，$f^{(n)}(x_0) > 0$时，$f$在$x_0$处不取极值。

综上，$n$为奇数时，$f$在$x_0$处不取极值。

\end{document}